% Adaptation du framework Style à un recueil global sous classe book.
% Version validée sur les 383 exercices du lot DYN.
\RequirePackage[a4paper,inner=22mm,outer=32mm,top=20mm,bottom=22mm,marginparwidth=24mm,marginparsep=4mm]{geometry}
\RequirePackage{multicol}
\RequirePackage{import}
\RequirePackage{float}
\RequirePackage{csquotes}
\RequirePackage{esvect}
\RequirePackage{cancel}
\RequirePackage{bm}
\RequirePackage{circuitikz}
\RequirePackage{fontawesome5}
\RequirePackage{fancyhdr}
\RequirePackage{titlesec}
\RequirePackage{marginnote}
\RequirePackage{caption}
\RequirePackage{comment}
\pgfplotsset{compat=1.18}

% Alias de couleurs historiques employés dans les environnements xpessoles.
\providecolor{White}{RGB}{255,255,255}
\providecolor{Black}{RGB}{0,0,0}
\providecolor{RoyalBlue}{RGB}{65,105,225}
\providecolor{ForestGreen}{RGB}{34,139,34}
\providecolor{Orange}{RGB}{255,165,0}
\providecolor{DarkOrange}{RGB}{255,140,0}
\providecolor{DarkKhaki}{RGB}{189,183,107}
\providecolor{Green}{RGB}{0,128,0}
\providecolor{ocre}{RGB}{204,119,34}
\providecolor{SIblue}{RGB}{31,78,121}
\providecolor{SIlightblue}{RGB}{229,240,250}
\providecolor{SIcorrigegreen}{RGB}{35,120,82}
\providecolor{SIcorrigelight}{RGB}{238,248,243}
\hypersetup{colorlinks=true,linkcolor=SIblue,urlcolor=SIblue,citecolor=SIblue}
\pagestyle{fancy}
\setlength{\headheight}{15pt}
\fancyhf{}
\lhead{Sciences industrielles pour l'ingénieur}
\rhead{Recueil d'exercices}
\cfoot{\thepage}
\titleformat{\chapter}[display]{\normalfont\huge\bfseries\color{SIblue}}{\chaptertitlename\ \thechapter}{10pt}{\Huge}
\titleformat{\section}{\Large\bfseries\color{SIblue}}{}{0pt}{}
\titleformat{\subsection}{\large\bfseries\color{SIblue}}{}{0pt}{}

\makeatletter
\@ifundefined{c@question}{\newcounter{question}}{}
\makeatother
\providecommand{\question}[1]{\refstepcounter{question}\par\medskip\noindent\textbf{Question \thequestion.}~#1\par\medskip}
\providecommand{\sidenote}[2][]{\marginnote{\footnotesize #2}}

% Notes marginales compactes et robustes pour plusieurs centaines d'exercices.
\renewcommand*{\marginfont}{\footnotesize}
\renewcommand{\marginnote}[2][]{\par\smallskip\noindent{\footnotesize\color{gray!70!black}#2}\par\smallskip}
\newenvironment{marginfigure}{\par\medskip\noindent\begin{minipage}{.82\linewidth}\centering\captionsetup{type=figure}}{\end{minipage}\par\medskip}

% -----------------------------------------------------------------------------
% Corrigés séparés : affichage global activable par une seule commande.
% -----------------------------------------------------------------------------
\newif\ifcorriges
\corrigesfalse
\newcommand{\AfficherCorriges}{\corrigestrue}
\newcommand{\MasquerCorriges}{\corrigesfalse}
\let\AvecCorriges\AfficherCorriges
\let\SansCorriges\MasquerCorriges

\newtcolorbox{corrigebox}[1][Corrigé]{
  enhanced,
  breakable,
  colback=SIcorrigelight,
  colframe=SIcorrigegreen,
  coltitle=white,
  fonttitle=\bfseries,
  title={#1},
  boxrule=.9pt,
  arc=1.5mm,
  left=3mm,right=3mm,top=2mm,bottom=2mm,
  before skip=5mm,after skip=5mm
}

\newcommand{\CorrigeQuestion}[1]{%
  \par\medskip\noindent{\bfseries\color{SIcorrigegreen}Question #1.}\enspace
}
\newcommand{\CorrigeACompleter}[1]{%
  \par\smallskip\noindent
  \textit{\color{gray!75!black}#1}%
  \par\smallskip
}
\newcommand{\CorrigeSource}[1]{%
  \par\smallskip\noindent{\scriptsize\color{gray!70!black}#1}\par\smallskip
}
\newcommand{\InclureCorrige}[1]{%
  \ifcorriges
    \par
    \IfFileExists{#1}{\input{#1}}{%
      \begin{corrigebox}[Corrigé indisponible]
        \CorrigeACompleter{Le fichier de corrigé attendu \texttt{#1} est absent.}
      \end{corrigebox}%
    }%
  \fi
}

% -----------------------------------------------------------------------------
% Compatibilité des blocs de code avec MiKTeX et l'encodage T1.
%
% environment_v2 contient une ancienne table listings « literate » qui remplace
% les caractères accentués par des commandes d'accent LaTeX. Certaines polices
% T1 de MiKTeX refusent ces commandes dès l'ouverture de lstlisting. On réinitialise
% réellement cette option (et non seulement sa valeur interne temporaire), puis on
% utilise une police machine à écrire en OT1, y compris pour verbatim.
% -----------------------------------------------------------------------------
\makeatletter
\@ifpackageloaded{listings}{%
  \lstset{
    literate={},
    inputencoding=utf8/latin1,
    upquote=false,
    basicstyle=\fontencoding{OT1}\ttfamily\small\selectfont
  }%
  \def\verbatim@font{\normalfont\fontencoding{OT1}\ttfamily\selectfont}%
}{}
\makeatother

% Les anciennes sources contiennent encore de nombreux blocs lstlisting dont
% l'initialisation échoue avec certaines installations MiKTeX/T1 avant même la
% lecture du code. Dans les recueils globaux, on ignore donc seulement ces blocs ;
% les sources Python restent présentes dans les fichiers individuels du dépôt.
\excludecomment{lstlisting}
